% =============================================================================
% Corrigés des exercices à faire seul
% =============================================================================
\chapter*{Corrigés des exercices à faire seul}
\addcontentsline{toc}{chapter}{Corrigés des exercices à faire seul}
\markboth{Corrigés des exercices à faire seul}{}

\epigraphe{Un corrigé ne sert qu'à celui qui a d'abord cherché.}

Tu trouveras ici le code, l'explication et les résultats de chaque exercice à faire seul. Les mêmes
corrigés, exécutables, sont dans le dossier \code{corriges/}. Si ton résultat diffère, ne corrige pas
ta requête pour « tomber juste » : cherche d'abord \emph{pourquoi} les deux diffèrent.

% =============================================================================
\section*{Chapitre 1}
\addcontentsline{toc}{section}{Chapitre 1}

\begin{correction}{1.3}
\textbf{Le journal.} L'ancien statut peut être NULL (premier statut d'un prêt). Avec
\code{CHECK (ancien_statut <> nouveau_statut)}, la ligne \code{(NULL, 'en_defaut')} serait acceptée,
mais par accident : la comparaison vaut NULL, et un \code{CHECK} laisse passer NULL.
\code{IS DISTINCT FROM} traite NULL comme une valeur ordinaire : \code{(NULL, 'en_defaut')} est acceptée
parce que les deux valeurs diffèrent vraiment, \code{('sain', 'sain')} est refusée. L'intention est
explicite, et la règle resterait juste si \code{nouveau_statut} devenait facultatif. Les valeurs par
défaut remplissent l'horodatage et l'utilisateur sans que les triggers aient à s'en occuper. Le journal
est consulté prêt par prêt : on indexe la clé étrangère.

\textbf{Les provisions.} La clé primaire composite \code{(date_arrete, pret_id)} garantit « une ligne par
prêt et par arrêté » ; elle sert aussi de cible à \code{ON CONFLICT} dans la clôture rejouable
(exercice 3.5). Toutes les colonnes sont \code{NOT NULL} : une provision incomplète n'a pas de sens.

\begin{lstlisting}[style=long]
-- 1. Journal d'audit : chaque changement de statut d'un prêt, avec la date et l'ancien statut.
CREATE TABLE IF NOT EXISTS journal_statuts (
    journal_id      BIGINT       GENERATED ALWAYS AS IDENTITY PRIMARY KEY,
    pret_id         BIGINT       NOT NULL REFERENCES demandes_pret (pret_id) ON DELETE RESTRICT,
    ancien_statut   VARCHAR(12),
    nouveau_statut  VARCHAR(12)  NOT NULL,
    jours_retard    INTEGER      CHECK (jours_retard >= 0),
    date_effet      DATE         NOT NULL,
    enregistre_le   TIMESTAMP    NOT NULL DEFAULT now(),
    enregistre_par  TEXT         NOT NULL DEFAULT session_user,   -- l'utilisateur connecté (voir chapitre 3, sécurité)
    CONSTRAINT statut_change CHECK (ancien_statut IS DISTINCT FROM nouveau_statut)
);
CREATE INDEX IF NOT EXISTS idx_journal_pret ON journal_statuts (pret_id);

-- 2. Provisions mensuelles : une photographie du risque de chaque prêt à chaque arrêté.
CREATE TABLE IF NOT EXISTS provisions_mensuelles (
    date_arrete   DATE          NOT NULL,
    pret_id       BIGINT        NOT NULL REFERENCES demandes_pret (pret_id) ON DELETE RESTRICT,
    jours_retard  INTEGER       NOT NULL CHECK (jours_retard >= 0),
    stage_ifrs9   SMALLINT      NOT NULL CHECK (stage_ifrs9 IN (1, 2, 3)),
    ead           NUMERIC(12,2) NOT NULL CHECK (ead >= 0),
    pd_utilisee   NUMERIC(7,6)  NOT NULL CHECK (pd_utilisee BETWEEN 0 AND 1),
    lgd           NUMERIC(4,3)  NOT NULL CHECK (lgd BETWEEN 0 AND 1),
    ecl           NUMERIC(12,2) NOT NULL CHECK (ecl >= 0),
    PRIMARY KEY (date_arrete, pret_id)       -- un prêt, une ligne par arrêté
);
\end{lstlisting}

\textbf{Pourquoi \code{RESTRICT} pour un journal ?} Un journal d'audit qui disparaîtrait avec le prêt
n'aurait plus aucune valeur de preuve. \textbf{Pourquoi pas de clé étrangère vers une table de
statuts ?} On aurait pu créer une petite table des statuts autorisés et la référencer ; ici, les
contraintes \code{CHECK} de \code{demandes_pret} garantissent déjà que seuls des statuts valides y
existent, et le journal ne fait que les recopier.
\end{correction}

\begin{correction}{1.4}
\textbf{a) Les règles par prêt.} On agrège chaque prêt dans une CTE, puis on compte les prêts en
anomalie. \code{IS DISTINCT FROM 0} attrape aussi un CRD final NULL (aucune échéance numéro $n$).
\begin{lstlisting}[style=long]
WITH par_pret AS (
    SELECT p.pret_id, p.montant_demande, p.duree_mois,
           sum(e.part_capital)                                                      AS capital_total,
           count(*)                                                                 AS nb_echeances,
           max(e.num_echeance)                                                      AS dernier_numero,
           max(e.capital_restant_du) FILTER (WHERE e.num_echeance = p.duree_mois)   AS crd_final,
           count(*) FILTER (WHERE e.montant_du <> e.part_capital + e.part_interets) AS montants_faux
    FROM demandes_pret p
    JOIN echeancier_remboursements e USING (pret_id)
    WHERE p.statut_dossier = 'accorde'
    GROUP BY p.pret_id
)
SELECT count(*)                                                                         AS prets_controles,
       count(*) FILTER (WHERE capital_total <> montant_demande)                           AS capital_different,
       count(*) FILTER (WHERE nb_echeances <> duree_mois OR dernier_numero <> duree_mois) AS echeances_manquantes,
       count(*) FILTER (WHERE crd_final IS DISTINCT FROM 0)                               AS crd_final_non_nul,
       sum(montants_faux)                                                                 AS montants_faux
FROM par_pret;
\end{lstlisting}

\textbf{b) Les règles entre deux échéances : l'auto-jointure.} La table apparaît deux fois, sous deux
alias : \code{prec} est l'échéance $k$, \code{suiv} l'échéance $k+1$ du même prêt. Grâce à l'index
unique \code{(pret_id, num_echeance)}, la jointure est rapide (0,3 seconde).
\begin{lstlisting}[style=long]
SELECT count(*) FILTER (WHERE suiv.capital_restant_du <> prec.capital_restant_du - suiv.part_capital) AS ruptures_crd,
       count(*) FILTER (WHERE suiv.date_echeance - prec.date_echeance NOT BETWEEN 28 AND 31)          AS ecarts_non_mensuels
FROM echeancier_remboursements prec
JOIN echeancier_remboursements suiv
  ON suiv.pret_id = prec.pret_id AND suiv.num_echeance = prec.num_echeance + 1;
\end{lstlisting}

\textbf{c)} Une simple jointure avec un filtre :
\begin{lstlisting}[style=long]
SELECT count(*) AS paiements_avant_octroi
FROM echeancier_remboursements e
JOIN demandes_pret p USING (pret_id)
WHERE e.date_paiement < p.date_decision;
\end{lstlisting}

\begin{resultat}
32\,410 prêts contrôlés ; \textbf{0} anomalie dans chacune des sept règles.
\end{resultat}

\textbf{Une histoire vraie.} La première version du corrigé de l'exercice 1.1 arrondissait séparément
les intérêts, le capital et le CRD. Ces contrôles ont trouvé \textbf{10 ruptures d'un centime}, toutes
sur le prêt d'Awa : le CRD d'un mois ne valait plus exactement celui du mois précédent moins la part de
capital. D'où la méthode actuelle (la part de capital est la différence de deux CRD arrondis). C'est
exactement à cela que servent des contrôles automatiques : trouver les erreurs qu'aucune relecture ne
voit.
\end{correction}

% =============================================================================
\section*{Chapitre 2}
\addcontentsline{toc}{section}{Chapitre 2}

\begin{correction}{2.5}
\begin{lstlisting}[style=long]
WITH test AS (
    SELECT ps.version_modele, ps.est_champion, ps.probabilite_defaut AS pd, ps.decision, s.en_defaut::int AS defaut
    FROM predictions_scoring_ml ps
    JOIN echantillons_modele em USING (pret_id)
    JOIN v_situation_prets s   USING (pret_id)
    WHERE em.echantillon = 'test'
),
rangs AS (
    -- ex æquo : RANK() donne le premier rang du groupe ; on prend le rang moyen du groupe
    SELECT *, RANK() OVER (PARTITION BY version_modele ORDER BY pd)
              + (count(*) OVER (PARTITION BY version_modele, pd) - 1) / 2.0 AS rang
    FROM test
),
mesures AS (
    SELECT version_modele, est_champion,
           (sum(rang) FILTER (WHERE defaut = 1) - sum(defaut) * (sum(defaut) + 1) / 2.0)
             / (sum(defaut) * sum(1 - defaut))                                          AS auc,
           100.0 * count(*) FILTER (WHERE decision = 'refuser' AND defaut = 1) / sum(defaut) AS rappel_pct,
           100.0 * count(*) FILTER (WHERE decision = 'refuser') / count(*)                AS taux_refus_pct,
           100.0 * count(*) FILTER (WHERE decision = 'accorder' AND defaut = 1)
                 / count(*) FILTER (WHERE decision = 'accorder')                          AS defaut_parmi_accordes_pct
    FROM rangs
    GROUP BY version_modele, est_champion
)
SELECT version_modele, est_champion,
       round(auc, 3)                     AS auc,
       round(2 * auc - 1, 3)             AS gini,
       round(rappel_pct, 1)              AS rappel_pct,
       round(taux_refus_pct, 1)          AS taux_refus_pct,
       round(defaut_parmi_accordes_pct, 1) AS defaut_parmi_accordes_pct
FROM mesures
ORDER BY auc DESC;
\end{lstlisting}

\textbf{Les points délicats.} Les deux fenêtres de \code{rangs} n'ont pas la même partition : la
première classe les prêts d'un modèle, la seconde compte les ex æquo d'une même PD. Les \code{2.0}
forcent des divisions décimales. \code{sum(1 - defaut)} compte les prêts sains.

\begin{resultat}
\begin{center}\small
\begin{tabular}{@{}lccccc@{}}
\toprule
Modèle & AUC & Gini & Rappel & Refus & Défaut si accordé \\
\midrule
\code{xgboost_monotone_v2} (champion) & \textbf{0,923} & \textbf{0,847} & 76,4~\% & \textbf{21,9~\%} & \textbf{6,6~\%} \\
\code{logistique_v1} & 0,876 & 0,751 & 77,8~\% & 29,5~\% & 6,9~\% \\
\code{grille_score_v1} & 0,875 & 0,750 & 77,9~\% & 29,6~\% & 6,9~\% \\
\bottomrule
\end{tabular}
\end{center}
\end{resultat}

\textbf{Le champion mérite-t-il sa place ?} Oui. Les challengers attrapent à peine plus de défauts
(77,8~\% contre 76,4~\%), mais en refusant \textbf{7,6 points de clients de plus} (29,5~\% contre
21,9~\%). Pour un risque un peu plus faible parmi les prêts accordés (6,6~\% contre 6,9~\%), le
XGBoost accepte donc beaucoup plus de bons clients. Le seuil de chaque modèle est un arbitrage métier
entre risque et volume : la base permet de le chiffrer.
\end{correction}

\begin{correction}{2.6}
\begin{lstlisting}[style=long]
WITH fins_de_mois AS (
    SELECT (d + INTERVAL '1 month - 1 day')::date AS date_arrete
    FROM generate_series(DATE '2025-09-01', DATE '2026-09-01', INTERVAL '1 month') AS d
),
-- Optimisation : ne garder que les échéances qui peuvent être impayées pendant la fenêtre
-- (échues avant la dernière fin de mois, pas encore soldées à la première). Filtrer AVANT de joindre : 7 s → 3 s.
impayes_possibles AS (
    SELECT pret_id, date_echeance, date_paiement
    FROM echeancier_remboursements
    WHERE date_echeance <= (SELECT max(date_arrete) FROM fins_de_mois)
      AND (date_paiement IS NULL OR date_paiement > (SELECT min(date_arrete) FROM fins_de_mois))
),
retards AS (
    SELECT f.date_arrete, p.pret_id,
           (p.date_premiere_echeance + (p.duree_mois - 1) * INTERVAL '1 month')::date AS fin_contrat,
           coalesce(f.date_arrete - min(e.date_echeance), 0)                           AS jours_retard
    FROM fins_de_mois f
    JOIN demandes_pret p ON p.statut_dossier = 'accorde' AND p.date_decision <= f.date_arrete
    -- LEFT JOIN : un prêt sans impayé n'a aucune ligne ici et garde un retard de 0
    LEFT JOIN impayes_possibles e
           ON e.pret_id = p.pret_id
          AND e.date_echeance <= f.date_arrete
          AND (e.date_paiement IS NULL OR e.date_paiement > f.date_arrete)
    GROUP BY f.date_arrete, p.pret_id
),
tranches AS (
    SELECT pret_id, date_arrete,
           CASE WHEN jours_retard = 0  THEN '1. À jour'
                WHEN jours_retard <= 30 THEN '2. 1 à 30 jours'
                WHEN jours_retard <= 60 THEN '3. 31 à 60 jours'
                WHEN jours_retard <= 90 THEN '4. 61 à 90 jours'
                ELSE                         '5. Plus de 90 jours' END AS tranche
    FROM retards
    WHERE jours_retard > 0 OR fin_contrat > date_arrete      -- on écarte les prêts terminés et soldés
),
transitions AS (
    SELECT tranche AS depart,
           LEAD(tranche) OVER (PARTITION BY pret_id ORDER BY date_arrete) AS arrivee   -- tranche le mois suivant
    FROM tranches
)
SELECT depart, count(*) AS nb_prets,
       round(100.0 * count(*) FILTER (WHERE arrivee = '1. À jour')           / count(*), 1) AS vers_a_jour,
       round(100.0 * count(*) FILTER (WHERE arrivee = '2. 1 à 30 jours')     / count(*), 1) AS vers_1_30,
       round(100.0 * count(*) FILTER (WHERE arrivee = '3. 31 à 60 jours')    / count(*), 1) AS vers_31_60,
       round(100.0 * count(*) FILTER (WHERE arrivee = '4. 61 à 90 jours')    / count(*), 1) AS vers_61_90,
       round(100.0 * count(*) FILTER (WHERE arrivee = '5. Plus de 90 jours') / count(*), 1) AS vers_plus_90
FROM transitions
WHERE arrivee IS NOT NULL                                     -- dernier mois observé : pas de mois suivant
GROUP BY depart
ORDER BY depart;
\end{lstlisting}

\textbf{La logique.} On fabrique un « panel » (prêt $\times$ fin de mois), on y calcule le retard
\emph{à cette date}, puis \code{LEAD} met côte à côte la tranche d'un mois et celle du mois suivant. Le
\code{WHERE arrivee IS NOT NULL} écarte le dernier mois, qui n'a pas de suivant.

\begin{resultat}[Ce que tu dois obtenir (en \% des prêts de la tranche de départ)]
\begin{center}\small
\newcommand{\hc}[1]{\cellcolor{rouge!#1}}
\begin{tabular}{@{}lrccccc@{}}
\toprule
Départ $\backslash$ mois suivant & Prêts & À jour & 1--30 & 31--60 & 61--90 & $> 90$ \\
\midrule
1. À jour & 177\,587 & \hc{40}93,8 & \hc{6}6,1 & 0,1 & 0,0 & 0,0 \\
2. 1 à 30 jours & 11\,674 & \hc{25}55,1 & \hc{3}5,3 & \hc{18}39,3 & 0,3 & 0,0 \\
3. 31 à 60 jours & 4\,862 & \hc{18}38,5 & 2,9 & 1,7 & \hc{25}55,9 & 1,0 \\
4. 61 à 90 jours & 2\,965 & \hc{4}7,7 & 1,1 & 0,3 & 0,7 & \hc{40}90,3 \\
5. Plus de 90 jours & 68\,213 & 0,0 & 0,0 & 0,0 & 0,0 & \hc{45}100,0 \\
\bottomrule
\end{tabular}
\end{center}
\end{resultat}

\textbf{Lecture.} Un prêt à jour le reste à 93,8~\%. Un prêt à 1--30 jours qui ne régularise pas dans
le mois passe mécaniquement à 31--60 (39,3~\%) : un mois de plus, c'est environ 30 jours de plus. Le
point décisif est la tranche 61--90 : \textbf{90,3~\%} de ces prêts finissent en défaut le mois suivant.
Le recouvrement doit donc agir dès 31--60 jours, quand 38,5~\% des prêts se régularisent encore. Enfin,
la tranche « plus de 90 jours » est absorbante (100~\%) : dans ces données, aucun prêt en défaut n'a
repris ses paiements.
\end{correction}

\begin{correction}{2.7}
\begin{lstlisting}[style=long]
WITH arrete AS (
    SELECT max(valeur_date) FILTER (WHERE cle = 'date_reference')     AS date_ref,
           max(valeur)      FILTER (WHERE cle = 'seuil_defaut_jours') AS seuil
    FROM parametres_risque
),
entrees_defaut AS (
    SELECT e.pret_id, min(e.date_echeance) + (a.seuil::int + 1) AS date_defaut
    FROM echeancier_remboursements e
    CROSS JOIN arrete a
    WHERE e.date_echeance <= a.date_ref
      AND coalesce(e.date_paiement, a.date_ref) - e.date_echeance > a.seuil
    GROUP BY e.pret_id, a.seuil
),
prets AS (
    SELECT p.pret_id, date_trunc('quarter', p.date_decision)::date AS generation,
           p.date_decision, d.date_defaut, a.date_ref
    FROM demandes_pret p
    CROSS JOIN arrete a
    LEFT JOIN entrees_defaut d USING (pret_id)
    WHERE p.statut_dossier = 'accorde'
),
courbes AS (
    SELECT g.generation, h.mois,
           count(*)                                                                    AS nb_observes,
           count(*) FILTER (WHERE g.date_defaut <= g.date_decision + h.mois * INTERVAL '1 month') AS nb_defauts
    FROM prets g
    CROSS JOIN (VALUES (6), (12), (18), (24)) AS h (mois)
    -- seuls les prêts déjà observés pendant h mois comptent, sinon les jeunes générations paraissent saines
    WHERE g.date_decision + h.mois * INTERVAL '1 month' <= g.date_ref
    GROUP BY g.generation, h.mois
),
tableau AS (
    SELECT generation,
           max(nb_observes) FILTER (WHERE mois = 6)                                   AS nb_prets,
           round(100.0 * sum(nb_defauts) FILTER (WHERE mois = 6)  / sum(nb_observes) FILTER (WHERE mois = 6), 1)  AS mob_6,
           round(100.0 * sum(nb_defauts) FILTER (WHERE mois = 12) / sum(nb_observes) FILTER (WHERE mois = 12), 1) AS mob_12,
           round(100.0 * sum(nb_defauts) FILTER (WHERE mois = 18) / sum(nb_observes) FILTER (WHERE mois = 18), 1) AS mob_18,
           round(100.0 * sum(nb_defauts) FILTER (WHERE mois = 24) / sum(nb_observes) FILTER (WHERE mois = 24), 1) AS mob_24
    FROM courbes
    GROUP BY generation
)
SELECT *, mob_12 - LAG(mob_12) OVER (ORDER BY generation) AS ecart_mob_12_vs_trimestre_precedent
FROM tableau
ORDER BY generation;
\end{lstlisting}

\begin{resultat}[Ce que tu dois obtenir (taux de défaut cumulé, en \%)]
\begin{center}\small
\begin{tabular}{@{}lrrrrrr@{}}
\toprule
Génération & Prêts & MOB 6 & MOB 12 & MOB 18 & MOB 24 & Écart MOB 12 \\
\midrule
T1 2022 & 1\,820 & 2,4 & 8,3 & 14,0 & 18,3 & \\
T2 2022 & 1\,894 & 2,3 & 7,7 & 11,9 & 16,5 & $-0{,}6$ \\
T3 2022 & 1\,922 & 1,8 & 8,5 & 13,0 & 17,1 & 0,8 \\
T4 2022 & 1\,836 & 1,4 & 9,0 & 13,5 & 18,4 & 0,5 \\
T1 2023 & 1\,904 & 2,3 & 9,7 & 14,5 & 18,2 & 0,7 \\
T2 2023 & 1\,900 & 2,8 & 9,4 & 15,4 & 19,4 & $-0{,}3$ \\
T3 2023 & 1\,867 & 1,8 & 9,0 & 14,1 & 18,7 & $-0{,}4$ \\
T4 2023 & 1\,879 & 2,0 & 9,8 & 14,6 & 20,2 & 0,8 \\
T1 2024 & 1\,883 & 2,5 & 10,4 & 15,9 & 20,6 & 0,6 \\
T2 2024 & 1\,843 & 3,3 & 10,8 & 18,0 & 23,3 & 0,4 \\
T3 2024 & 1\,839 & 2,2 & 10,6 & 17,9 & 24,6 & $-0{,}2$ \\
T4 2024 & 1\,931 & 2,0 & 12,4 & 19,0 & & 1,8 \\
T1 2025 & 1\,819 & 4,3 & 14,3 & 23,5 & & 1,9 \\
\rowcolor{rougeFond} T2 2025 & 1\,840 & 5,0 & 17,6 & & & 3,3 \\
\rowcolor{rougeFond} T3 2025 & 1\,825 & 5,1 & 23,1 & & & 5,5 \\
T4 2025 & 1\,568 & 2,1 & & & & \\
T1 2026 & 1\,409 & 0,0 & & & & \\
\bottomrule
\end{tabular}
\end{center}
\end{resultat}

\textbf{Lecture, et une leçon importante.} Les générations 2022--2023 sont stables (MOB 12 entre 8 et
10~\%). À partir de fin 2024, le MOB 12 s'envole jusqu'à 23,1~\%. Dans une vraie banque, ce serait un
signal d'alerte majeur : la politique d'octroi se serait relâchée. \textbf{Ici, c'est un artefact de la
simulation} : chaque prêt qui a réellement fait défaut devait être en défaut au 30/09/2026 ; pour les
prêts récents, cela oblige le défaut à arriver tôt. Pour la même raison, la génération T1 2026 n'a aucun
défaut : les dates des prêts en défaut ont été tirées avant novembre 2025. Réflexe à garder :
\emph{avant d'interpréter un signal, se demander s'il vient du phénomène étudié ou de la façon dont les
données ont été construites.}
\end{correction}

\begin{correction}{2.8}
\textbf{a)} On classe dans une CTE (avec la part de chaque prêt dans l'encours en défaut de sa note,
grâce à une seconde fenêtre), puis on filtre le rang dans la requête englobante.
\begin{lstlisting}[style=long]
WITH classement AS (
    SELECT s.note_risque, s.pret_id, c.nom_complet, s.encours, s.jours_retard,
           RANK() OVER (PARTITION BY s.note_risque ORDER BY s.encours DESC) AS rang,
           round(100 * s.encours / sum(s.encours) OVER (PARTITION BY s.note_risque), 2) AS part_de_la_note_pct
    FROM v_situation_prets s
    JOIN clients c USING (client_id)
    WHERE s.en_defaut
)
SELECT * FROM classement WHERE rang <= 3 ORDER BY note_risque, rang;
\end{lstlisting}

\begin{resultat}[Extrait (21 lignes au total, noms fictifs)]
\begin{center}\small
\begin{tabular}{@{}clrrrcr@{}}
\toprule
Note & Client & Prêt & Encours (€) & Retard & Rang & Part de la note \\
\midrule
A & Tristan Fischer & 13 & 32\,580,44 & 446 & 1 & 0,46~\% \\
A & Adrienne Dias & 29141 & 31\,397,09 & 243 & 2 & 0,44~\% \\
A & Daniel Pruvost & 32159 & 31\,257,48 & 1\,052 & 3 & 0,44~\% \\
\multicolumn{7}{c}{\dots} \\
G & Antoine Lopez & 14759 & 23\,836,89 & 470 & 1 & 2,92~\% \\
G & Henri Colin & 12221 & 23\,692,79 & 180 & 2 & 2,91~\% \\
G & Thomas Robert de la Renard & 23159 & 23\,469,55 & 245 & 3 & 2,88~\% \\
\bottomrule
\end{tabular}
\end{center}
\end{resultat}

\textbf{b) Le signal précoce.} Deux \code{LAG} sur la même fenêtre nommée, et \code{ROW_NUMBER}
décroissant pour repérer la dernière échéance payée.
\begin{lstlisting}[style=long]
WITH delais AS (
    SELECT pret_id, num_echeance,
           date_paiement - date_echeance                                              AS retard,
           LAG(date_paiement - date_echeance, 1) OVER w                               AS retard_precedent,
           LAG(date_paiement - date_echeance, 2) OVER w                               AS retard_avant_precedent,
           ROW_NUMBER() OVER (PARTITION BY pret_id ORDER BY num_echeance DESC)        AS recence
    FROM echeancier_remboursements
    WHERE date_paiement IS NOT NULL
    WINDOW w AS (PARTITION BY pret_id ORDER BY num_echeance)
)
SELECT d.pret_id, s.note_risque, d.retard_avant_precedent, d.retard_precedent, d.retard, s.jours_retard, s.encours
FROM delais d
JOIN v_situation_prets s USING (pret_id)
WHERE d.recence = 1                                            -- la dernière échéance payée
  AND d.retard > d.retard_precedent AND d.retard_precedent > d.retard_avant_precedent
  AND d.retard > 0
  AND s.jours_retard <= 30
  AND s.encours > 0                                            -- piège : sans ce filtre, on surveille des prêts soldés
ORDER BY d.retard DESC, s.encours DESC
LIMIT 20;
\end{lstlisting}

\begin{resultat}[Les cinq premiers des 20 prêts signalés]
\begin{center}\small
\begin{tabular}{@{}rcrrrrr@{}}
\toprule
Prêt & Note & Retard $k-2$ & Retard $k-1$ & Retard $k$ & Retard actuel & Encours (€) \\
\midrule
27484 & A & $-2$ & $-1$ & 40 & 26 & 11\,230,97 \\
23306 & A & $-1$ & 37 & 40 & 26 & 1\,759,94 \\
31436 & C & $-3$ & $-2$ & 31 & 13 & 11\,167,63 \\
15739 & C & $-2$ & 0 & 20 & 0 & 32\,651,38 \\
27395 & C & $-2$ & $-1$ & 20 & 0 & 21\,543,39 \\
\bottomrule
\end{tabular}
\end{center}
\end{resultat}

\textbf{Pourquoi le filtre de l'ordre d'exécution ?} \code{recence}, \code{retard_precedent}\dots{} sont
des résultats de fenêtres : impossible de les filtrer dans le \code{WHERE} de la CTE qui les calcule.
On les filtre dans la requête suivante. \textbf{Usage métier :} ces prêts ne sont pas encore en
difficulté, mais leurs paiements dérivent ; un appel du conseiller coûte moins cher qu'un recouvrement.
\end{correction}

% =============================================================================
\section*{Chapitre 3}
\addcontentsline{toc}{section}{Chapitre 3}

\begin{correction}{3.4}
\textbf{a)} La vérité terrain est chargée dans une table temporaire, qui disparaîtra à la fin de la
session. La requête compte, sur les mêmes prêts, les défauts réels, les défauts détectés et les
désaccords (\code{<>} entre deux booléens = « ou exclusif »).
\begin{lstlisting}[style=long]
CREATE TEMP TABLE verite_defauts (pret_id BIGINT PRIMARY KEY, defaut_reel SMALLINT NOT NULL);
\copy verite_defauts FROM 'donnees/verite_defauts.csv' WITH (FORMAT csv, HEADER true)
SELECT count(*) FILTER (WHERE v.defaut_reel = 1)                                   AS defauts_reels,
       count(*) FILTER (WHERE d.statut_pret = 'en_defaut')                         AS defauts_detectes,
       count(*) FILTER (WHERE (d.statut_pret = 'en_defaut') <> (v.defaut_reel = 1)) AS desaccords   -- attendu : 0
FROM demandes_pret d
JOIN verite_defauts v USING (pret_id);
\end{lstlisting}

\textbf{b)} Grâce à la date de référence paramétrée, « faire passer le temps » est une simple mise à
jour, annulée à la fin.
\begin{lstlisting}[style=long]
BEGIN;
UPDATE parametres_risque SET valeur_date = DATE '2027-01-31' WHERE cle = 'date_reference';
SELECT statut_pret AS avant_traitement FROM demandes_pret WHERE pret_id = 32582;      -- « sain » : aucun événement
CALL actualiser_statuts(32582);
SELECT statut_pret AS apres_traitement, fn_jours_retard(32582, fn_date_reference()) AS jours_retard
FROM demandes_pret WHERE pret_id = 32582;                                            -- « en_defaut », 98 jours
SELECT nouveau_statut, date_effet FROM journal_statuts WHERE pret_id = 32582;
ROLLBACK;
\end{lstlisting}

\begin{resultat}
a) 7\,088 défauts réels, 7\,088 détectés, \textbf{0 désaccord}.\par
b) Avant le traitement : \code{sain} (aucun événement ne s'est produit). Après : \code{en_defaut}, avec
\textbf{98 jours} de retard. Le journal contient la ligne \code{en_defaut}, date d'effet 31/01/2027.
\end{resultat}

\textbf{Ce que prouve b).} Le passage du temps seul n'a rien déclenché : le statut était encore « sain »
avec 98 jours de retard. C'est le traitement planifié qui l'a corrigé, et le trigger de journal a alors
tracé le changement et envoyé l'alerte. Les deux mécanismes sont complémentaires.
\end{correction}

\begin{correction}{3.5}
\begin{lstlisting}[style=long]
CREATE OR REPLACE PROCEDURE cloture_mensuelle(p_date_arrete DATE)
LANGUAGE plpgsql AS $$
DECLARE
    v_nb  INTEGER;
    v_lgd NUMERIC := fn_parametre('lgd');
BEGIN
    IF p_date_arrete > fn_date_reference() THEN
        RAISE EXCEPTION 'Arrêté du % impossible : postérieur à la date de référence (%)',
                        p_date_arrete, fn_date_reference();
    END IF;
    INSERT INTO provisions_mensuelles (date_arrete, pret_id, jours_retard, stage_ifrs9, ead, pd_utilisee, lgd, ecl)
    SELECT p_date_arrete, s.pret_id, s.jours_retard, c.stage, s.encours, c.pd, v_lgd,
           round(c.pd * v_lgd * s.encours, 2)          -- l'ECL se recalcule à partir des colonnes stockées
    FROM fn_situation_prets(p_date_arrete) s
    JOIN demandes_pret d          ON d.pret_id = s.pret_id
    JOIN predictions_scoring_ml ps ON ps.pret_id = s.pret_id AND ps.est_champion
    CROSS JOIN LATERAL (SELECT fn_stage_ifrs9(s.jours_retard) AS stage) st
    CROSS JOIN LATERAL (SELECT st.stage,
                               round(fn_pd_ifrs9(st.stage, ps.probabilite_defaut, s.mois_restants, d.duree_mois), 6) AS pd) c
    WHERE s.encours > 0
    ON CONFLICT (date_arrete, pret_id) DO UPDATE           -- rejouer un arrêté remplace, ne duplique pas
        SET jours_retard = EXCLUDED.jours_retard, stage_ifrs9 = EXCLUDED.stage_ifrs9, ead = EXCLUDED.ead,
            pd_utilisee  = EXCLUDED.pd_utilisee,  lgd = EXCLUDED.lgd, ecl = EXCLUDED.ecl;
    GET DIAGNOSTICS v_nb = ROW_COUNT;
    RAISE NOTICE 'Arrêté du % : % prêts provisionnés', p_date_arrete, v_nb;
END $$;

-- Historique : les 12 derniers arrêtés de fin de mois (un CALL ne s'utilise pas dans un SELECT : bloc DO)
DO $$
DECLARE
    v_mois DATE;
BEGIN
    FOR v_mois IN
        SELECT (d + INTERVAL '1 month - 1 day')::date
        FROM generate_series(DATE '2025-10-01', DATE '2026-09-01', INTERVAL '1 month') AS d
    LOOP
        CALL cloture_mensuelle(v_mois);
    END LOOP;
END $$;

-- Évolution des provisions et dotation nette du mois (ce qui pèse sur le résultat de la banque)
WITH mensuel AS (
    SELECT date_arrete,
           sum(ead)                                     AS ead,
           sum(ecl)                                     AS ecl,
           sum(ecl) FILTER (WHERE stage_ifrs9 = 1)      AS ecl_stage1,
           sum(ecl) FILTER (WHERE stage_ifrs9 = 2)      AS ecl_stage2,
           sum(ecl) FILTER (WHERE stage_ifrs9 = 3)      AS ecl_stage3
    FROM provisions_mensuelles
    GROUP BY date_arrete
)
SELECT date_arrete, ead, ecl, ecl_stage1, ecl_stage2, ecl_stage3,
       ecl - LAG(ecl) OVER (ORDER BY date_arrete)       AS dotation_nette,
       round(100 * ecl / ead, 2)                        AS taux_couverture_pct
FROM mensuel
ORDER BY date_arrete;

-- Contrôle : rejouer le dernier arrêté ne change rien (idempotence)
CALL cloture_mensuelle(DATE '2026-09-30');
SELECT count(*) AS lignes_au_30_09, sum(ecl) AS ecl_au_30_09
FROM provisions_mensuelles WHERE date_arrete = DATE '2026-09-30';
\end{lstlisting}

\textbf{\code{CROSS JOIN LATERAL}.} Une sous-requête \code{LATERAL} peut utiliser les colonnes des tables
qui la précèdent : la première calcule le stage, la seconde la PD qui dépend du stage. C'est une façon
lisible d'enchaîner des calculs ligne par ligne sans répéter les expressions.

\begin{resultat}[Ce que tu dois obtenir (montants en millions d'euros)]
\begin{center}\small
\begin{tabular}{@{}lrrrrrrr@{}}
\toprule
Arrêté & EAD & ECL & Stage 1 & Stage 2 & Stage 3 & Dotation nette & Couverture \\
\midrule
31/10/2025 & 147,01 & 21,54 & 5,03 & 1,35 & 15,16 & & 14,66~\% \\
30/11/2025 & 146,64 & 21,94 & 4,57 & 1,39 & 15,97 & $+0{,}39$ & 14,96~\% \\
31/12/2025 & 146,48 & 22,32 & 3,95 & 1,44 & 16,92 & $+0{,}38$ & 15,24~\% \\
31/01/2026 & 146,26 & 22,71 & 3,40 & 1,51 & 17,80 & $+0{,}40$ & 15,53~\% \\
28/02/2026 & 145,93 & 23,13 & 2,87 & 1,52 & 18,75 & $+0{,}42$ & 15,85~\% \\
31/03/2026 & 145,84 & 23,62 & 2,31 & 1,57 & 19,74 & $+0{,}48$ & 16,19~\% \\
30/04/2026 & 145,79 & 24,17 & 1,83 & 1,46 & 20,88 & $+0{,}55$ & 16,58~\% \\
31/05/2026 & 145,56 & 24,71 & 1,42 & 1,32 & 21,97 & $+0{,}54$ & 16,97~\% \\
30/06/2026 & 145,62 & 25,15 & 1,02 & 1,22 & 22,92 & $+0{,}44$ & 17,27~\% \\
31/07/2026 & 141,94 & 25,46 & 0,96 & 0,59 & 23,91 & $+0{,}31$ & 17,93~\% \\
31/08/2026 & 137,79 & 25,67 & 0,91 & 0,03 & 24,74 & $+0{,}22$ & 18,63~\% \\
30/09/2026 & 133,96 & 25,63 & 0,86 & 0,03 & 24,74 & $-0{,}05$ & 19,13~\% \\
\bottomrule
\end{tabular}
\end{center}
Après avoir rejoué le 30/09 : 21\,501 lignes, ECL 25\,627\,130,87~€, inchangés.
\end{resultat}

\textbf{Lecture.} Chaque mois, de nouveaux prêts passent en stage 3 : la provision augmente, et la
\textbf{dotation nette} (environ 0,4 million par mois) pèse sur le résultat de la banque. La provision
de stage 1 fond de 5,03 à 0,86 million : les prêts qui allaient faire défaut avaient une PD élevée et
pesaient lourd en stage 1 ; en passant en stage 3, ils laissent en stage 1 des prêts peu risqués. En
septembre, plus aucun nouveau défaut (le stage 3 ne bouge plus) : la dotation devient négative, c'est
une \textbf{reprise de provision}.

\textbf{Un écart de 1,72~€ ?} Le total au 30/09 (25\,627\,130,87~€) diffère de celui de l'exercice 2.2
(25\,627\,129,15~€). Ici, la PD est arrondie à 6 décimales et l'ECL au centime \emph{prêt par prêt},
avant la somme. Savoir expliquer un écart d'arrondi, c'est montrer qu'on contrôle ses chiffres.
\end{correction}

\begin{correction}{3.6}
\begin{lstlisting}[style=long]
DO $$
BEGIN
    IF NOT EXISTS (SELECT FROM pg_roles WHERE rolname = 'analyste_risque') THEN
        CREATE ROLE analyste_risque NOLOGIN;
    END IF;
    IF NOT EXISTS (SELECT FROM pg_roles WHERE rolname = 'application_paiement') THEN
        CREATE ROLE application_paiement NOLOGIN;
    END IF;
END $$;

-- L'analyste lit tout, ne modifie rien
GRANT SELECT ON ALL TABLES IN SCHEMA public TO analyste_risque;

-- L'application d'encaissement n'écrit PAS dans les tables : elle appelle la procédure, qui s'exécute
-- avec les droits de son propriétaire (SECURITY DEFINER). search_path figé : sinon, un utilisateur
-- pourrait créer une fausse table ou fonction homonyme et détourner la procédure.
ALTER PROCEDURE enregistrer_paiement(BIGINT, NUMERIC, DATE) SECURITY DEFINER SET search_path = public, pg_temp;
REVOKE ALL ON PROCEDURE enregistrer_paiement(BIGINT, NUMERIC, DATE) FROM PUBLIC;
GRANT EXECUTE ON PROCEDURE enregistrer_paiement(BIGINT, NUMERIC, DATE) TO application_paiement;

-- Vérification (transaction annulée)
BEGIN;
SET LOCAL ROLE application_paiement;
CALL enregistrer_paiement(32582, 100);                          -- autorisé : passe par la procédure
DO $$
BEGIN
    UPDATE echeancier_remboursements SET montant_paye = montant_du WHERE pret_id = 32582;
    RAISE NOTICE 'ANOMALIE : écriture directe acceptée';
EXCEPTION WHEN insufficient_privilege THEN
    RAISE NOTICE 'Écriture directe refusée, comme prévu';
END $$;
ROLLBACK;
\end{lstlisting}

\begin{resultat}
\code{NOTICE: Prêt 32582 : 100 € encaissés le 2026-09-30}, puis
\code{NOTICE: Écriture directe refusée, comme prévu}.
\end{resultat}

\textbf{Pourquoi \code{NOLOGIN} ?} Ce sont des \emph{groupes} de droits. Les vrais comptes (l'analyste
Paul, le service d'encaissement) seraient des rôles \code{LOGIN} membres de ces groupes
(\code{GRANT analyste_risque TO paul}). \code{SET LOCAL ROLE} ne dure que jusqu'à la fin de la
transaction : le \code{ROLLBACK} rend aussi ton rôle d'origine.
\end{correction}
